\section{The Names, the Letters, and Number}

Kabbalah builds the world out of \textbf{letters} and \textbf{numbers}. The
\textbf{Book of Formation} (Sefer Yetzirah) teaches that God made creation by
combining the twenty-two Hebrew letters. The divine \textbf{names} are the next
layer up, whole structures built by spelling and summing those letters. And
\textbf{gematria}, the reading of letters as numbers, is the tool that ties words
and meanings together.

\subsection{The Twenty-Two Letters as Building Blocks}

The Sefer Yetzirah counts \textbf{thirty-two paths of wisdom}. These are the
\textbf{ten sefirot} plus the \textbf{twenty-two letters} (10 + 22 = 32). The
twenty-two letters split into three classes by their cosmic role.

\begin{itemize}
 \item \textbf{Three mothers}. The elemental letters, tied to the three
 elements.
 \item \textbf{Seven doubles}. Each has a hard and a soft sound. Tied to the
 seven planets, the seven days, and the six directions plus center.
 \item \textbf{Twelve simples}. Each has a single sound. Tied to the twelve
 zodiac signs, the twelve months, and the twelve tribes.
\end{itemize}

Three plus seven plus twelve makes twenty-two. The table below gives the three
classes and their chief correspondences.

\begin{longtable}{L{0.16\linewidth} L{0.26\linewidth} L{0.46\linewidth}}
\caption{The twenty-two Hebrew letters in their three classes.}\\
\toprule
\textbf{Class} & \textbf{Letters} & \textbf{Maps to} \\
\midrule
Three mothers & aleph, mem, shin & air, water, fire (the three elements) \\
Seven doubles & bet, gimel, dalet, kaf, pe, resh, tav & seven planets, seven days, six directions and center \\
Twelve simples & he, vav, zayin, chet, tet, yod, lamed, nun, samech, ayin, tzadi, qof & twelve zodiac signs, twelve months, twelve tribes \\
\bottomrule
\end{longtable}

The three mothers carry the elements. \textbf{Aleph} is air and breath,
\textbf{mem} is water, \textbf{shin} is fire. They sit at the heart of the
letter-cosmology, the silent, the mute, and the hissing sounds.

\subsection{The Sealing of the Six Directions}

Before the world was set, God sealed the six directions of space, each with a
permutation of the three letters \textbf{YHV} (yod, he, vav, the first three
letters of the great name). Three distinct letters give exactly six orderings,
one per direction. The six directions plus the center make seven, which is why
they tie to the seven double letters.

\subsection{The Divine Names}

The root divine name is the \textbf{Tetragrammaton}, written \textbf{YHVH} (yod,
he, vav, he). Its plain value is \textbf{26}. From it, Lurianic Kabbalah builds a
whole architecture by the \textbf{filling} method (milui), writing out the full
name of each letter and summing the result. The four classic fillings are named by
their totals. They are called \textbf{Av}, \textbf{Sag}, \textbf{Mah}, and
\textbf{Ban}.

\begin{longtable}{L{0.12\linewidth} L{0.12\linewidth} L{0.30\linewidth} L{0.34\linewidth}}
\caption{The four expansions of the Tetragrammaton.}\\
\toprule
\textbf{Name} & \textbf{Value} & \textbf{Stage} & \textbf{Maps to} \\
\midrule
Av & 72 & wisdom, the seminal point & Abba (father), world of Atzilut \\
Sag & 63 & understanding, the vessels that shattered & Imma (mother), Atzilut to Beriah \\
Mah & 45 & repair, the new name of rectification & Zeir Anpin (the son), Beriah to Yetzirah \\
Ban & 52 & kingship, the receiving vessel & Nukva (the daughter), Yetzirah to Asiyah \\
\bottomrule
\end{longtable}

Two of these carry special weight. \textbf{Sag} (63) is tied to the
\textbf{breaking of the vessels}, the world of points whose vessels could not hold
the light. \textbf{Mah} (45) is the name of \textbf{rectification}, and it is also
the value of \textbf{Adam} (man, 45). So Mah is the name of the repairing human
force. The higher pair Av and Sag are the parent names. The lower pair Mah and Ban
are the repairing and receiving names.

Beyond the four-letter name and its fillings, the tradition counts a family of
longer names, each a fixed string of letters with its own use.

\begin{longtable}{L{0.20\linewidth} L{0.58\linewidth}}
\caption{The principal name-constructions by length.}\\
\toprule
\textbf{Name} & \textbf{What it is} \\
\midrule
The 4-letter name & YHVH itself, value 26. Its twelve distinct permutations map to the twelve months and tribes. \\
The 42-letter name & Encoded in the prayer Ana BeKoach, a seven-line poem. Tied to the act of creation and the ascent through the worlds. \\
The 72-letter name & The Shem HaMephorash, seventy-two three-letter triplets woven from three verses of Exodus, each verse holding seventy-two letters. \\
\bottomrule
\end{longtable}

The seventy-two-letter name is the largest of the family. It is drawn from Exodus
14, verses 19, 20, and 21, each of which has exactly seventy-two letters. Writing
the three verses in a grid and reading down the columns yields seventy-two
three-letter names. In later tradition each is paired with an angel and a degree
of the zodiac and used in meditation.

\subsection{The Thirteen Attributes of Mercy}

Beyond the spelled-out fillings, the tradition holds a second great name-formula,
the \textbf{thirteen attributes of mercy} (Shelosh Esreh Middot). It is not a
spelling of YHVH but a fixed verse-formula, drawn from \textbf{Exodus 34}, verses
6 and 7, the words God speaks to Moses after the golden calf. The thirteen
attributes are recited in penitential prayer to draw down mercy. In Lurianic
Kabbalah they are the inner content of the \textbf{thirteen tikkunim of the beard}
of Arikh Anpin, the all-merciful face.

\begin{longtable}{L{0.06\linewidth} L{0.26\linewidth} L{0.52\linewidth}}
\caption{The thirteen attributes of mercy from Exodus 34.}\\
\toprule
\textbf{No.} & \textbf{Attribute} & \textbf{Sense} \\
\midrule
1 & The name (YHVH) & Mercy before a person sins. \\
2 & The name repeated (YHVH) & Mercy after a person sins and repents. \\
3 & God (El) & The aspect of mercy and power. \\
4 & Compassionate (Rachum) & Merciful, sparing the suffering. \\
5 & Gracious (Chanun) & Giving freely even to the undeserving. \\
6 & Slow to anger (Erekh Apayim) & Patient, holding back judgment. \\
7 & Abounding in kindness (Rav Chesed) & Great in loving-kindness. \\
8 & And truth (Emet) & Faithful and reliable in mercy. \\
9 & Keeping kindness for thousands (Notzer Chesed La-Alafim) & Preserving kindness across generations. \\
10 & Bearing iniquity (Nose Avon) & Forgiving deliberate wrongdoing. \\
11 & And transgression (Nose Pesha) & Forgiving rebellion. \\
12 & And sin (Nose Chata'ah) & Forgiving error and the unintended. \\
13 & Cleansing (Nakeh) & Wiping the slate of those who repent. \\
\bottomrule
\end{longtable}

The verse opens with the divine name written \textbf{twice}, YHVH YHVH. The two
namings are counted as the first two attributes. They are read as \textbf{mercy
before a person sins and mercy after}. The exact partition of the thirteen varies
between authorities, but the count of thirteen is fixed, and the formula always
runs from the doubled name through the final cleansing.

\subsection{Gematria and the Methods}

\textbf{Gematria} reads the numerical value of letters and words. Two words with
the same value are held to share a hidden link. The first nine letters are the
units, the next nine the tens, and the last four the hundreds. So aleph is 1, yod
is 10, qof is 100, tav is 400. A few widely cited equalities show the method at
work.

\begin{itemize}
 \item \textbf{Echad} (one) equals \textbf{Ahavah} (love), both \textbf{13}.
 Oneness and love share a value.
 \item \textbf{Elohim} (God) equals \textbf{HaTeva} (nature), both \textbf{86}.
 God and nature share a value.
 \item \textbf{Adam} (man) equals \textbf{45}, the value of the name Mah, the
 force of repair.
\end{itemize}

Beyond plain counting, three classic letter operations open the hidden sense of
scripture.

\begin{itemize}
 \item \textbf{Notarikon}. Reading a word as an acronym, or compressing a phrase
 into one word from its initial letters. The word PaRDeS is itself a
 notarikon for the four levels of reading.
 \item \textbf{Temurah}. Letter substitution by a fixed table. The best-known
 tables are \textbf{Atbash}, which swaps the first letter with the last and
 so on inward, and \textbf{Albam}, which splits the alphabet in half and
 pairs the halves.
 \item \textbf{Gematria} proper, the value-summing above, which links words by
 number.
\end{itemize}

These tools serve the deepest of the four levels of reading. The four levels are
named by the word \textbf{PaRDeS}, an orchard.

\begin{longtable}{L{0.10\linewidth} L{0.20\linewidth} L{0.48\linewidth}}
\caption{The four levels of reading scripture (PaRDeS).}\\
\toprule
\textbf{Letter} & \textbf{Level} & \textbf{Meaning} \\
\midrule
P & Peshat & the plain, literal sense \\
R & Remez & the hint, the allegorical sense \\
D & Derash & the homiletic, seeking and comparing sense \\
S & Sod & the secret, the mystical or Kabbalistic sense \\
\bottomrule
\end{longtable}

\textbf{The secret} (Sod) is the level Kabbalah works in. The letter-methods are
the keys to Sod. The four levels are sometimes mapped onto the four worlds, Peshat
to the lowest and Sod to the highest, so that reading deeper is climbing the same
ladder the cosmos descends.
